\chapter{LiDAR and Obstacle Guard}

\section{YDLIDAR role}

The YDLIDAR is a USB serial sensor. Its ROS 2 driver publishes \topic{/scan} as \cmd{sensor_msgs/msg/LaserScan}. The scan is used by RViz, the obstacle guard, and future SLAM/navigation layers. The LiDAR is not wired into the Sabertooth chain; it influences motion only through the ROS bridge.

The current known-good LiDAR configuration is in:

\begin{lstlisting}[language=bash]
src/rudra_base_bridge/config/ydlidar_g2b.yaml
\end{lstlisting}

Important parameters include:

\begin{center}
\small
\begin{tabularx}{\linewidth}{p{0.22\linewidth}Y}
\toprule
Parameter & Value\\
\midrule
\param{port} & \filep{/dev/serial/by-id/usb-Silicon_Labs_CP2102_USB_to_UART_Bridge_Controller_0001-if00-port0}\\
\param{baudrate} & 128000\\
\param{frame_id} & \cmd{laser}\\
\param{sample_rate} & 5\\
\param{frequency} & 10.0\\
\param{range_min} & 0.28 m\\
\param{range_max} & 16.0 m\\
\param{intensity} & true\\
\bottomrule
\end{tabularx}
\end{center}

\section{LaserScan geometry}

A LaserScan message contains an array of ranges. Each index maps to an angle:

\begin{equation}
\theta_i=\theta_{min}+i\Delta\theta.
\end{equation}

Each valid point can be converted into LiDAR-frame Cartesian coordinates:

\begin{align}
x_i &= r_i\cos\theta_i,\\
y_i &= r_i\sin\theta_i.
\end{align}

The RUDRA obstacle guard does not build a map. It computes only the closest valid range inside a front cone.

\begin{figure}[htbp]
\centering
\resizebox{0.98\linewidth}{!}{\begin{tikzpicture}[scale=2.1]
\draw[->,thick] (0,0) -- (2.2,0) node[right]{front x};
\draw[->,thick] (0,-1.2) -- (0,1.2) node[above]{y};
\draw[rudraTeal,thick] (0,0) -- ({2*cos(35)},{2*sin(35)});
\draw[rudraTeal,thick] (0,0) -- ({2*cos(-35)},{2*sin(-35)});
\draw[rudraTealSoft,fill=rudraTealSoft,opacity=0.55] (0,0) -- ({2*cos(35)},{2*sin(35)}) arc (35:-35:2) -- cycle;
\draw[rudraAmber,thick,dashed] (0,0) circle (0.45);
\draw[rudraTeal,thick,dashed] (0,0) circle (1.0);
\node at (1.15,0.62) {front cone};
\node[rudraAmber] at (0.62,-0.18) {stop};
\node[rudraTeal] at (1.26,-0.18) {slow};
\fill[rudraAmber] (0.55,0.2) circle (0.04) node[above right]{closest obstacle};
\end{tikzpicture}}
\caption{Obstacle guard geometry. Only finite scan points inside the configured front cone are considered.}
\end{figure}

\section{Guard scale law}

The current guard parameters are:

\begin{align}
\phi_{front} &= 70^\circ,\\
d_{stop} &= 0.45\ \text{m},\\
d_{slow} &= 1.00\ \text{m},\\
t_{scan} &= 0.75\ \text{s}.
\end{align}

Let \(d\) be the closest valid front-cone obstacle distance. The forward scale is:

\begin{equation}
s(d)=
\begin{cases}
1, & \text{guard disabled},\\
1, & \text{no fresh scan},\\
1, & d \ge d_{slow},\\
0, & d \le d_{stop},\\
\dfrac{d-d_{stop}}{d_{slow}-d_{stop}}, & d_{stop}<d<d_{slow}.
\end{cases}
\end{equation}

For positive forward commands, the filtered command is:

\begin{equation}
u_{forward,filtered}=s(d)u_{forward}.
\end{equation}

Reverse and turn-in-place remain available so the robot can back away from an obstacle.

\begin{figure}[htbp]
\centering
\resizebox{0.98\linewidth}{!}{\begin{tikzpicture}
\begin{axis}[width=0.82\linewidth,height=6cm,xlabel={closest front distance (m)},ylabel={forward scale},xmin=0,xmax=1.4,ymin=-0.05,ymax=1.05,grid=both]
\addplot[thick,rudraTeal,domain=0:0.45,samples=2]{0};
\addplot[thick,rudraTeal,domain=0.45:1.0,samples=50]{(x-0.45)/(1.0-0.45)};
\addplot[thick,rudraTeal,domain=1.0:1.4,samples=2]{1};
\addplot[thick,rudraAmber,dashed] coordinates {(0.45,-0.05) (0.45,1.05)};
\addplot[thick,black,dashed] coordinates {(1.0,-0.05) (1.0,1.05)};
\node at (axis cs:0.23,0.12) {blocked};
\node at (axis cs:0.72,0.55) {slowing};
\node at (axis cs:1.18,0.92) {clear};
\end{axis}
\end{tikzpicture}}
\caption{Forward-speed scaling used by the obstacle guard.}
\end{figure}

\section{Marker and status output}

The guard publishes:

\begin{itemize}
  \item \topic{/rudra/obstacle_guard}: string status containing reason, enabled flag, closest front distance, and forward scale.
  \item \topic{/rudra/obstacle_guard_marker}: RViz marker at the closest front obstacle point.
\end{itemize}

Marker colors are:

\begin{center}
\begin{tabular}{ll}
\toprule
Marker color & Meaning\\
\midrule
Green & clear\\
Orange & slowing\\
Red & blocked\\
Gray & disabled\\
Deleted/no marker & no closest front point or stale scan\\
\bottomrule
\end{tabular}
\end{center}

\section{LiDAR startup behavior}

The YDLIDAR driver may take 10--15 seconds before the first valid scan appears. Startup checksum warnings and repeated fixed-point warnings can occur while scans are still valid. The operational test is not a warning-free console; it is whether \topic{/scan} publishes with frame \cmd{laser} and the scan is visible in RViz.

\section{Launching LiDAR}

LiDAR only:

\begin{lstlisting}[language=bash]
cd /home/rudra/Projects/RUDRAv4
source /opt/ros/lyrical/setup.bash
source /home/rudra/ros2_ws/install/setup.bash
source install/setup.bash
bash scripts/run_lidar.sh \
  /dev/serial/by-id/usb-Silicon_Labs_CP2102_USB_to_UART_Bridge_Controller_0001-if00-port0
\end{lstlisting}

View from Aghora:

\begin{lstlisting}[language=bash]
source /opt/ros/lyrical/setup.bash
export ROS_DOMAIN_ID=42
export ROS_LOCALHOST_ONLY=0
cd /home/rudra/Projects/RUDRAv4
source install/setup.bash
ros2 launch rudra_base_bridge lidar_view.launch.py
\end{lstlisting}

The viewer launch starts RViz only. It does not start the LiDAR driver and does not start a drive bridge.
